\section{Walking on the Edge of an Abyss}

When people approach spiritual ideas, they typically have preconceived notions of what an enlightened being or master should be like. Perhaps they expect someone to dress a certain way, speak in a sentimental manner, hold liberal political views, or demonstrate great powers over the material world. These notions will either attract them to the legion of charismatic gurus who know how to play the con game, or cause them to reject those nondescript types, who really do have a higher level of understanding. Evola addresses this topic in the final chapter of \textbf{The Hermetic Tradition}.


There are two things to keep in mind.

\begin{enumerate}


\item As Evola points out, the magus or initiate is not at all interested in manifesting unusual phenomena in order to astonish, amaze or terrify. From his perspective, there is simply no point to it.

\item If an ordinary person already knows how an initiate should act and comport himself, then what sets the two apart? The magus has an understanding of the world that the uninitiated do not, and that needs be reflected in many ways.
\end{enumerate}


Ordinary people bristle when they hear this, because they are egalitarians at heart, either explicitly or implicitly. Overt egalitarians on the left consider the initiate to be nothing more than an idealized version of themselves; hence, rather than being challenged by the initiate, they instead find confirmation in their own beliefs. Thus their egalitarianism is protected.

Those on the so-called right reject out-of-hand any claims to occult knowledge. They may claim to rely on the authority of reason, and trust in science, materialism, genetics or so on. The alternative is to accept a worldview based on faith in place of a genuine gnosis or knowledge. Evola addresses the latter two options in ``Those who know and those who believe'' from \textit{Pagan Imperialism}\footnote{\url{https://www.amazon.com/Pagan-Imperialism-Julius-Evola/dp/0999086006/ref=sr\_1\_1?crid=4U5MSAZZ3RCE&dchild=1&keywords=pagan+imperialism+evola&qid=1626944181&s=books&sprefix=evola+pagan\%2Caps\%2C195&sr=1-1}}.

In opposition to preconceived notions, the initiate will appear quite ordinary. The difference lies in their interiority, which the uninitiated are incapable of recognizing. Evola offers this portrayal of the initiate.

\begin{quotex}
The initiate is an occult being and his path is neither visible nor penetrable. He is elusive, not to be pigeonholed. He arrives from the direction contrary to that towards which all gazes are fixed and takes the most natural seeming vehicle for his supernatural action. He may be an intimate friend, companion, or lover; he may be sure of possessing all your heart and confidence. But he will always be something different, other than what he lets be known. We will perceive this ``other'' only when we have penetrated his domain. And then perhaps we will have the feeling of having been \textbf{walking on the edge of an abyss}.
\end{quotex}



\flrightit{Posted on 2010-07-22 by Cologero}
